\section{Introduction}

This chapter describes the main threats to validity and limitations of the PatchOps thesis. The goal is not to
weaken the contribution, but to state precisely where the evidence applies. PatchOps is an MTech-scale research
prototype evaluated on bounded datasets and controlled harnesses. Its strongest claims concern safety-gated
decision-making, abstention, localization bottlenecks, and measured negative results. Its weaker claims concern
general repair success, broad production readiness, and semantic correctness beyond the available oracles.

The limitations are organized using four standard categories: construct validity, internal validity, external
validity, and conclusion validity. The chapter then summarizes the practical engineering limitations of the
current system.

\section{Construct Validity}

Construct validity asks whether the thesis measures the concepts it claims to measure. In PatchOps, this issue
is most visible around repair success, localization quality, safety, and patch plausibility.

\subsection{CI Pass Is Not Full Correctness}

The thesis uses CI Pass@1 because it is the natural benchmark signal for CI repair. However, a passing CI run is
not equivalent to semantic correctness. A patch may pass the current tests while violating intended behavior,
weakening assertions, deleting checks, or changing only the symptom. This limitation is fundamental to CI repair
and APR, not only to PatchOps.

PatchOps mitigates this threat by explicitly separating CI validation from semantic correctness. The thesis
uses reward-hacking detection and patch-plausibility analysis to avoid treating every green or syntactically
valid patch as trustworthy. Even so, these defenses are partial. They can catch specific dishonest or off-target
patterns, but they cannot prove that a patch is functionally equivalent to the developer's intended fix.

\subsection{File-Overlap Is a Proxy for On-Target Repair}

The plausibility analysis uses overlap with the developer's changed files as a proxy for whether a generated
patch is on target. This is useful because it is objective and scalable, but it is imperfect. A correct patch
may edit a different file from the developer fix, especially if there are multiple legitimate repair locations.
Conversely, a patch may edit a gold-overlapping file while still being semantically wrong.

For this reason, the thesis describes plausibility as an on-target signal rather than as a correctness oracle.
The result supports a safer auto-PR decision, not a proof of repair correctness.

\subsection{Localization Metrics Capture Ranking, Not Understanding}

Recall@k and Hit@k measure whether the true fix file appears in a ranked candidate list. These metrics are
appropriate for comparing localizers, but they do not fully capture repository understanding. A localizer may
rank the right file because of a lexical coincidence, or it may understand the failure but rank a nearby file
outside the gold set. This affects all localization methods in the thesis.

The mitigation is comparative rather than absolute. PatchOps evaluates lexical, LLM, embedding, and agentic
approaches under the same gold-file criterion. The conclusion that these methods plateau below the target bar is
therefore meaningful for the measured setup, even if the metric is not a complete measure of reasoning.

\subsection{Safety Labels Are Operational}

Safety in this thesis is defined operationally: forbidden file classes, secret or permission risks, deployment
risks, flaky rerun cases, workflow weakening, reward-hacking patterns, and inappropriate automatic patching.
This makes the safety benchmark measurable, but it does not cover every possible safety issue in a real
repository. For example, a patch could be business-logic unsafe without touching an obviously forbidden file.

The thesis therefore claims measured safety under defined policies, not universal safety. This is still useful:
CI agents need explicit policies, and PatchOps demonstrates how such policies can be integrated into the repair
decision.

\section{Internal Validity}

Internal validity asks whether the observed results are caused by the intended experimental factors rather than
by bugs, leakage, or uncontrolled variation.

\subsection{Pipeline Bugs and Regression Risk}

PatchOps is a multi-stage system. A bug in log extraction, repository cloning, ranking, diff generation, or
decision routing could affect multiple reported metrics. The project mitigates this risk through regression
tests, acceptance tests, per-phase artifacts, and explicit amendments when bugs are discovered. For example, the
localization ordering issue was corrected before the later localization conclusions were frozen.

Nevertheless, implementation bugs remain possible. The thesis should be read as evidence from the current
artifact, not as a proof that every pipeline component is bug-free.

\subsection{Model Non-Determinism}

LLM calls are not perfectly reproducible. Even with low-temperature settings and fixed prompts, provider-side
model updates, decoding differences, and transient API behavior can change outputs. This risk is visible in the
agentic localization phase, where abstention behavior differed across runs.

PatchOps mitigates this risk by logging phase results, using deterministic components where possible, and
adding a learned local classifier for the classification stage. The limitation remains for LLM-based
localization and patch generation. Future repeated-run studies would strengthen variance estimates.

\subsection{Local CI Reproduction}

The Docker validation harness replays a restricted subset of GitHub Actions behavior. It is not a full GitHub
Actions emulator. Differences in hosted runners, caches, service containers, custom actions, secrets, matrix
jobs, and operating-system behavior can affect whether a failure reproduces locally.

The thesis handles this by reporting reproducibility separately from repair success. Non-reproducible cases are
not silently counted as repaired or failed repairs. However, the repair results apply only to the reproducible
subset under the local harness.

\subsection{Budget and Provider Constraints}

Some experiments use cost-bounded model calls. This is appropriate for a practical thesis, but it also affects
which methods are explored deeply. The agentic localizer, for example, was evaluated under a bounded step and
budget regime with mid-tier and limited stronger-model probes. A larger budget, stronger reasoning model, or
more sophisticated tool interface could change the absolute localization numbers.

The thesis therefore frames the agentic result as an honest negative for the measured configuration, not as a
universal impossibility result for all future agentic localization.

\section{External Validity}

External validity asks whether the results generalize beyond the evaluated datasets, languages, and repositories.

\subsection{Python-First Benchmark Scope}

The comparable benchmark results are Python-first because CI-Repair-Bench is Python-based. PatchOps also has
live demonstration cases outside that exact benchmark scope, but those are treated as demonstrations rather
than as benchmark evidence. The results should not be generalized automatically to JavaScript, Java, Go, Rust,
mobile builds, or polyglot monorepos.

The main architectural ideas--safety gates, abstention, CI validation, plausibility checks, and reward-hack
detection--are language-agnostic in principle. The measured effectiveness of the current implementation is not.

\subsection{GitHub Actions Scope}

PatchOps targets GitHub Actions. Other CI systems such as GitLab CI, Jenkins, CircleCI, Buildkite, and Azure
Pipelines have different configuration languages, permission models, caching behavior, log formats, and failure
patterns. The thesis findings may transfer conceptually, but the implementation and taxonomy would need
adaptation.

\subsection{Benchmark Representativeness}

CI-Repair-Bench provides real failures and is an appropriate benchmark for this project, but no benchmark can
cover the full distribution of CI failures in the wild. The public benchmark is also low-risk dominated compared
with the kinds of secret, permission, release, and deployment failures where abstention policy matters most.
PatchOps addresses this by adding a labeled safety set, but that set is smaller and partly synthetic.

As a result, the repair and localization results have stronger external validity for public Python CI failures
than for high-risk enterprise pipelines. The safety results demonstrate policy behavior under defined cases, but
larger real-world safety datasets would be needed for stronger generalization.

\subsection{Synthetic-Heavy Decision-Gate Data}

The Phase 15 learned decision gate uses the expanded 76-case Safety Set. This directly improves coverage over
the original 20-case safety set, but it also introduces an important validity threat: most of the expanded cases
are hand-authored synthetic cases. Such data is useful for testing policy behavior under labeled risk
conditions, including adversarial reward-hack temptations, but it may also encode the author's labeling
shortcuts. The Phase 15 result should therefore be read as an indicative labeled-set result, not as a field
measurement from deployed CI pipelines.

\subsection{Repository Size and Complexity}

The experiments include real repositories, but repository complexity varies widely in practice. Very large
monorepos, generated code, vendored dependencies, custom build systems, and service-backed integration tests may
change both localization and validation behavior. The current retrieval and agentic tools are bounded and may
not scale linearly to such settings.

\section{Conclusion Validity}

Conclusion validity asks whether the evidence is strong enough to support the stated conclusions.

\subsection{Sample Sizes}

Several experiments have modest sample sizes. The scaled diagnosis ablation uses \(n=30\), the repair linkage
uses \(n=24\) reproducible cases, the original Phase 5 safety result uses \(n=20\), the Phase 15 decision-gate
study uses \(n=76\), and the patch audit uses 66 produced patches. The held-out classifier and
selective-prediction split is larger at \(n=118\), but it still does not eliminate sampling uncertainty.

The thesis mitigates this by avoiding over-precise claims and by reporting negative results conservatively. For
example, it does not claim that agentic localization can never work; it claims that the measured bounded agent
does not beat the baselines on the evaluated slice. It does not claim zero universal risk; it reports zero
unsafe auto-fixes under the evaluated safety definitions.

For the learned decision gate specifically, the Wilson confidence intervals overlap with the hand-tuned
threshold baseline at \(n=76\). The thesis therefore uses the phrase ``matches-or-beats'' and emphasizes removal
of the hand-set thresholds rather than claiming a statistically separated victory.

\subsection{Multiple Experiments and Narrative Selection}

Because PatchOps went through many phases, there is a risk of selecting only the most favorable results. The
writing mitigates this by foregrounding negative results and preserving the phase chronology in the source
documents. The thesis narrative is built around the measured bottlenecks, not around a cherry-picked success
case.

\subsection{Threshold Choices}

Several gates depend on thresholds: classifier confidence, plausibility score, risk level, and decision routing.
Thresholds are necessary for a deployable system, but the exact defaults may not be optimal. The thesis treats
them as operating points on a coverage-risk curve rather than as universal constants. Phase 15 reduces this
threat for the low-risk act-versus-hold decision by replacing the old hand-set threshold with a learned
calibrated gate, but the learned gate still has its own operating point and inherits the limitations of the
Safety Set used to train and evaluate it.

\section{Practical System Limitations}

PatchOps is not production-ready in the broad sense. It demonstrates an end-to-end architecture and measured
research claims, but several engineering limitations remain.

First, the local validation harness supports only a restricted subset of common Python CI workflows. A complete
system would need stronger handling for matrix builds, service containers, custom actions, caches, secrets,
container jobs, and operating-system differences.

Second, the current repair generator is not a strong semantic repair engine. The experiments show that once
patches apply, semantic correctness remains the wall. Future work needs stronger program understanding,
differential testing, invariant checking, or developer-in-the-loop review.

Third, the repository navigator is bounded and model-limited. It provides useful evidence, but it is not yet a
robust replacement for deterministic retrieval. Better tools, graph or AST structure, and stronger reasoning
models may improve it.

Fourth, the safety policy is necessarily incomplete. It captures important CI-specific hazards, but each
organization would need to adapt forbidden paths, protected workflows, deployment rules, and approval policies.

\section{Chapter Conclusion}

The main limitations of PatchOps are bounded scope, modest sample sizes, proxy correctness metrics, local CI
reproduction limits, and incomplete generalization beyond Python GitHub Actions. These limitations do not
invalidate the central contribution, because the thesis is careful about what it claims: PatchOps is a measured
safety-gated and abstention-aware CI repair system, not a universal CI repair solver. The next chapter concludes
the thesis and outlines future work that follows directly from these limits.
